\documentclass[11pt]{article}

\usepackage[margin=1in]{geometry}
\usepackage{amsmath}
\usepackage{amssymb}
\usepackage{booktabs}
\usepackage{tabularx}
\usepackage{hyperref}

\title{Run Parameters for the NIFTy8 Healpix Denoising Runs}
\author{}
\date{}

\begin{document}
\maketitle

\section{Inference Problem}

The synthetic and real-data scripts use the same inference problem.  For a
Healpix map with $N_{\rm pix}=12N_{\rm side}^2$ pixels,
\begin{equation}
    d = s + n,
    \qquad
    n_p \sim \mathcal{N}(0,\sigma_n^2),
\end{equation}
with independent white noise in pixel space.  The posterior is
\begin{equation}
    p(s,\theta \mid d)
    \propto
    \exp\left[
        -\frac{1}{2\sigma_n^2}
        \sum_{p=1}^{N_{\rm pix}} (d_p - s_p)^2
    \right]
    p(s\mid\theta)p(\theta),
\end{equation}
where $\theta$ denotes the latent prior/hyperprior variables used by NIFTy.

NIFTy samples this posterior approximately with variational inference.  The
run parameters below control the numerical approximation, not the mathematical
prior model.

\section{Resolution and Data Controls}

\begin{center}
\small
\begin{tabularx}{\textwidth}{lXX}
\toprule
Parameter & Mathematical role & Practical effect \\
\midrule
\texttt{NSIDE}
& Sets $N_{\rm pix}=12N_{\rm side}^2$
& Resolution of the input map. \\
\texttt{ANALYSIS\_NSIDE}
& Sets the $N_{\rm side}$ actually inferred on
& Optional downsampled run for testing. \\
\texttt{COMPONENT}
& Selects one scalar field $d$
& Which I/Q/U FITS component is read. \\
\texttt{NOISE\_SIGMA}
& Sets $\sigma_n$ in the likelihood
& Controls the weight of data residuals. \\
\texttt{SYNTH\_SNR}
& Sets $\sigma_n = {\rm std}(s)/{\rm SNR}$
& Synthetic runs only; changes generated noise and likelihood together. \\
\bottomrule
\end{tabularx}
\end{center}

\paragraph{Important.}
If \texttt{ANALYSIS\_NSIDE} differs from \texttt{NSIDE}, the map is degraded
before inference.  The original pixel-white noise is then no longer exactly
white and independent at the degraded resolution.  Such runs are useful as
sanity checks, but the full statistical model is cleanest when
\texttt{ANALYSIS\_NSIDE = NSIDE}.

\section{Sampling and Variational Inference Controls}

NIFTy minimizes a variational approximation to the posterior.  Let $m_i$ denote
the current latent mean at outer VI iteration $i$.  At a high level, each
outer iteration performs
\begin{align}
    \{\xi_{i,k}\}_{k=1}^{K}
    &\sim q_i(\xi\mid m_i),
    \\
    m_{i+1}
    &\approx
    \arg\min_m
    \widehat{\mathrm{KL}}_i(m),
\end{align}
where $K=\texttt{POSTERIOR\_SAMPLES}$ and
$\widehat{\mathrm{KL}}_i$ is a sample-based approximation to the KL objective.

\begin{center}
\small
\begin{tabularx}{\textwidth}{lXX}
\toprule
Parameter & Mathematical role & Practical effect \\
\midrule
\texttt{POSTERIOR\_SAMPLES}
& Number $K$ of samples used per VI iteration
& More samples reduce Monte Carlo noise but cost more. \\
\texttt{SAVE\_SAMPLES}
& Number of final signal samples written
& Controls output size, not the inference objective. \\
\texttt{VI\_ITERATIONS}
& Number of outer VI refinements
& More rounds improve the posterior approximation but multiply cost. \\
\texttt{GEOVI}
& Chooses geoVI if true, MGVI-like mode if false
& geoVI can be more accurate but more expensive. \\
\texttt{GEOVI\_ITERATIONS}
& Nonlinear solve cap inside geoVI sampling
& Used only when \texttt{GEOVI=True}. \\
\texttt{GEOVI\_TOL}
& Nonlinear geoVI stopping tolerance
& Used only when \texttt{GEOVI=True}. \\
\bottomrule
\end{tabularx}
\end{center}

\paragraph{\texttt{VI\_ITERATIONS}.}
This is the number of outer posterior-refinement cycles.  It is not the number
of saved samples.  For example,
\[
    \texttt{POSTERIOR\_SAMPLES}=10,
    \qquad
    \texttt{VI\_ITERATIONS}=3
\]
means that NIFTy uses 10 approximate samples in each of 3 outer VI updates.
The final saved maps are drawn from the final approximation.

\section{Optimization Controls}

Inside each VI iteration, NIFTy solves optimization and metric-inversion
subproblems.  These are controlled by iteration caps and tolerances.

\begin{center}
\small
\begin{tabularx}{\textwidth}{lXX}
\toprule
Parameter & Mathematical role & Practical effect \\
\midrule
\texttt{MAP\_ITERATIONS}
& Iteration cap for minimizing $\widehat{\mathrm{KL}}$
& Larger values allow better mean/hyperparameter updates. \\
\texttt{MAP\_TOL}
& Gradient-norm stopping threshold
& Smaller values demand tighter optimization. \\
\texttt{SAMPLE\_CG\_ITERATIONS}
& Iteration cap for metric inverse solves
& Larger values improve posterior sample accuracy. \\
\texttt{SAMPLE\_CG\_TOL}
& CG stopping threshold
& Smaller values demand more accurate samples. \\
\bottomrule
\end{tabularx}
\end{center}

When NIFTy prints
\begin{quote}
\texttt{Iteration limit reached. Assuming convergence}
\end{quote}
one of these inner iteration caps has been reached.  The run may still produce
usable outputs, but the diagnostics should then be interpreted as approximate.
Increasing \texttt{MAP\_ITERATIONS} or \texttt{SAMPLE\_CG\_ITERATIONS} usually
improves convergence at increased runtime.

\section{Randomness and Output Controls}

\begin{center}
\small
\begin{tabularx}{\textwidth}{lXX}
\toprule
Parameter & Mathematical role & Practical effect \\
\midrule
\texttt{SEED}
& Initializes pseudo-random draws
& Controls reproducibility. \\
\texttt{OUTDIR}
& None
& Directory for FITS samples, metadata, and progress logs. \\
\texttt{DRY\_RUN}
& None
& Builds/checks setup and exits before inference. \\
\bottomrule
\end{tabularx}
\end{center}

\section{Cost Scaling}

The number of pixels scales as
\begin{equation}
    N_{\rm pix} = 12N_{\rm side}^2.
\end{equation}
Going from $N_{\rm side}=512$ to $1024$ multiplies the pixel count by 4.
Because each VI step requires repeated spherical harmonic transforms and
linear solves, the wall time can increase by more than this factor.

A conservative production sequence is:
\begin{align*}
    &N_{\rm side}=128 \quad \text{synthetic test},\\
    &N_{\rm side}=256 \quad \text{real-data sanity test},\\
    &N_{\rm side}=512 \quad \text{intermediate real-data test},\\
    &N_{\rm side}=1024 \quad \text{full run}.
\end{align*}

\end{document}
